\begin{pastexam}[label={px:ntut-110-ci-03}]{北科大資工所 110 學年度 第四題}
（問答）

Suppose Alice and Bob work on a project at different places and Bob needs to send a binary tree $T$
to Alice. The communication channel between Bob and Alice only allows to send/receive one node of
the tree at a time. One way for them to communicate is to serialize the binary tree. A binary tree
can be serialized by arranging the nodes linearly. Thus, a \textit{serialization} of a given binary
tree $T$ is a sequence of nodes. Answer the following questions. (10 pts)
\begin{enumerate}[label=(\arabic*)]
\item How many serializations for a given binary tree $T$ with $n$ nodes? (2 pts)
\item Please provide a way for Bob to send a binary tree $T$ over the channel and Alice can restore
$T$ by receiving the nodes from the channel. Explain how Bob sends and Alice restores the tree with
time and space analysis. (4 pts)
\item Suppose that the tree to be sent is a \textit{binary search tree} (BST). Again, please provide
a more efficient way for them to communication and explain why the scheme works with time and space
analysis. (4 pts)
\end{enumerate}
\end{pastexam}

\begin{pastsoln}
題目對「序列化」的定義只說是「把節點排成一列」，沒有要求能還原。 (1) 依這個定義作答；(2)(3) 要
設計的是「能還原的」序列化，下面假設每個節點送出時可以附帶少量資訊（例如兩個位元），鍵值互不相同。

\textbf{(1) $n!$ 種。} 序列化就是 $n$ 個（不同的）節點的一種排列，第一個位置有 $n$ 種選法、第二個有
$n - 1$ 種……，共 $n!$ 種。 其中只有少數排法配上適當的資訊才能還原樹，這就是 (2) 要處理的事。

\textbf{(2) 送前序，每個節點附帶「有沒有左子、有沒有右子」兩個位元。}

Bob 照前序（根、左子樹、右子樹）一個一個送；Alice 照同樣的順序收，收到一個節點就看它的兩個位元，
決定要不要接著遞迴收它的左子樹、右子樹。

\begin{pseudo}{Bob 送出、Alice 還原（前序加兩個位元）}
\Procedure{Send}{$x$}
    \State 送出 $(x.\mathit{key},\ x.\mathit{left} \ne \NIL,\ x.\mathit{right} \ne \NIL)$
    \If{$x.\mathit{left} \ne \NIL$}
        \State \Call{Send}{$x.\mathit{left}$}
    \EndIf
    \If{$x.\mathit{right} \ne \NIL$}
        \State \Call{Send}{$x.\mathit{right}$}
    \EndIf
\EndProcedure
\Statex
\Function{Receive}{}
    \State 收下一個節點 $(k, \mathit{hasL}, \mathit{hasR})$，建立新節點 $x$，$x.\mathit{key} \gets k$
    \State $x.\mathit{left} \gets$ \textbf{if} $\mathit{hasL}$ \textbf{then} \Call{Receive}{} \textbf{else} \NIL
    \State $x.\mathit{right} \gets$ \textbf{if} $\mathit{hasR}$ \textbf{then} \Call{Receive}{} \textbf{else} \NIL
    \State \Return $x$
\EndFunction
\end{pseudo}
Bob 呼叫 \proc{Send}($T.\mathit{root}$)（$T$ 非空），Alice 呼叫 \proc{Receive}() 得到根。

正確性：對子樹的節點數做歸納。 \proc{Send}($x$) 送出的正是以 $x$ 為根的子樹的前序，而且連續送完才
輪到別的節點。 假設 Alice 呼叫 \proc{Receive} 時，通道上接下來是某棵子樹 $S$ 的前序：第一個收到
的是 $S$ 的根，它的兩個位元正確說出左右子樹是否存在；若左子樹存在，接下來連續的就是左子樹的前序，
由歸納假設，遞迴呼叫恰好把它還原、而且剛好收完；右子樹同理。 所以 \proc{Receive} 還原出 $S$，
左右不會弄反。 這其實就是資結書第~7.4~節的延伸前序（範例~7.1 用 \texttt{.} 表示空的子樹）：把「後面
跟一個空的子樹」的資訊放進節點本身的位元，就不用另外送 $n + 1$ 個空標記。 如果通道只能送真正的
節點、不能附帶位元，也可以送前序和中序兩串（共 $2n$ 次），用資結書虛擬碼~7.5 還原。

時間：Bob 每個節點送一次，$\Theta(n)$；Alice 每收一個節點做 $O(1)$ 的事，$\Theta(n)$。 通訊量是
$n$ 個節點加 $2n$ 個位元。 空間：Bob 只用遞迴堆疊，$O(h)$，$h$ 是樹高（最壞 $n$）；Alice 要存整棵樹
$\Theta(n)$，再加遞迴堆疊 $O(h)$。

\textbf{(3) 二元搜尋樹：只送前序的鍵，不用附帶任何位元。}

Bob 照前序送出鍵值。 Alice 用「每棵子樹的鍵一定落在某個範圍內」來判斷收到的鍵屬於哪裡：

\begin{pseudo}{由二元搜尋樹的前序還原}
\Function{Receive-BST}{$\mathit{lo}, \mathit{hi}$}
    \If{通道已經沒有鍵 \OR 下一個鍵 $k$ 不在 $(\mathit{lo}, \mathit{hi})$ 之間}
        \State \Return \NIL \Comment{只偷看，不取出}
    \EndIf
    \State 取出 $k$，建立新節點 $x$，$x.\mathit{key} \gets k$
    \State $x.\mathit{left} \gets \Call{Receive-BST}{\mathit{lo}, k}$
    \State $x.\mathit{right} \gets \Call{Receive-BST}{k, \mathit{hi}}$
    \State \Return $x$
\EndFunction
\end{pseudo}
Alice 呼叫 \proc{Receive-BST}($-\infty, +\infty$)。 「偷看下一個鍵」只要 Alice 多收一個、先放在一個
變數裡就好。

為什麼可行：二元搜尋樹的前序裡，根之後先是左子樹的所有鍵（都小於根），再來是右子樹的所有鍵（都大於
根）（資結書第~9.1~節）。 以節點 $x$ 為根的子樹，它的鍵都落在一個開區間 $(\mathit{lo}, \mathit{hi})$
裡，這兩個界就是從根走到 $x$ 的路上，最近一次往右轉、往左轉的那兩個祖先的鍵。 呼叫
\proc{Receive-BST}($\mathit{lo}, \mathit{hi}$) 時，若下一個鍵在範圍內，它一定是這個位置的子樹的根；
若不在範圍內，表示這個位置是空的，下一個鍵屬於某個祖先的右子樹（比 $\mathit{hi}$ 大）。 對子樹大小
歸納，每個呼叫恰好取走它那棵子樹的鍵，還原出來的樹和 Bob 的一樣。

時間：每一次 \proc{Receive-BST} 不是取走一個鍵（共 $n$ 次），就是直接回傳 \NIL（每個空指標一次，共
$n + 1$ 次），每次 $O(1)$，總共 $\Theta(n)$。 通訊量只有 $n$ 個鍵，比 (2) 少了 $2n$ 個位元。 空間：
Alice 存樹 $\Theta(n)$，遞迴堆疊 $O(h)$；Bob 做前序走訪，$O(h)$。

比較：把收到的鍵依序用一般的插入（資結書虛擬碼~9.6）插進空的二元搜尋樹，也會得到同一棵樹（前序裡
父節點總是比子孫先到），但每次插入要從根走下去，總時間 $O(nh)$，樹退化成一條鏈時是 $\Theta(n^2)$。
用範圍判斷的做法每個鍵只處理一次，是線性的。
\end{pastsoln}
